% ---- score2 standard preamble (from: python3 score2_tools.py preamble) ----
\documentclass[10pt]{extarticle}
\usepackage[a4paper,margin=1.5cm]{geometry}
\usepackage[T1]{fontenc}
\usepackage{amsmath,amssymb}
\usepackage{graphicx}
\usepackage[export]{adjustbox}
\usepackage{multicol}
\usepackage{enumitem}
\usepackage{booktabs,array,tabularx}
\setlength{\columnsep}{0.7cm}
\setlength{\parindent}{0pt}
\setlength{\parskip}{3pt}
\pagestyle{empty}
\makeatletter
% figure registry (figdefs.tex is generated by the extract step)
\newcommand{\qfigdef}[4]{\@namedef{qf@file@#1}{#2}\@namedef{qf@w@#1}{#3}\@namedef{qf@wide@#1}{#4}}
\newcommand{\qf@check}[1]{\@ifundefined{qf@file@#1}{\PackageError{score2}{Unknown figure id #1}{}}{}}
\newcommand{\qf@img}[4][]{\includegraphics[width=#2\dimexpr\@nameuse{qf@w@#4}pt\relax,max width=#3,max height=0.45\textheight,#1]{\@nameuse{qf@file@#4}}}
% Full-width questions: finalize wraps every question that contains a wide figure in \qfullbegin...\qfullend,
% so the WHOLE question (stem, figures, options) is set across the page in reading order, then the two columns
% resume. The question's head - its number, stem and first figure (\qfullsplit marks the end of the head) -
% is kept on one page, so a number is never left at a page foot with its figure overleaf; the rest may break
% across pages normally, so no big gaps are left. Without a \qfullsplit, the whole question is kept together
% if it is under 45% of a page. Never write these yourself.
\newif\ifqf@full
\newif\ifqf@boxed
\newsavebox{\qf@box}
\newcommand{\qf@open}{\global\qf@fulltrue\par\global\qf@boxedtrue
  \setbox\qf@box\vbox\bgroup\hsize\textwidth\linewidth\textwidth}
\newcommand{\qf@emit}[1]{\par\egroup\global\qf@boxedfalse
  \ifdim\dimexpr\ht\qf@box+\dp\qf@box\relax<#1\textheight\noindent\usebox{\qf@box}\par\else\unvbox\qf@box\fi}
\newcommand{\qfullsplit}{\ifqf@boxed\qf@emit{0.85}\fi}
\newcommand{\qf@close}{\ifqf@boxed\qf@emit{0.45}\fi\global\qf@fullfalse}
\newcommand{\qfullbegin}{\end{multicols}\qf@open}
\newcommand{\qfullend}{\qf@close\begin{multicols}{2}}
\newcommand{\qfullstart}{\qf@open}
\newcommand{\qfulldone}{\qf@close}
% \qfig[scale]{id}: block figure, centred, directly where it is written (size set by the tool).
% Fallback if a wide figure is not inside a full-width question: break out of the columns in place - no floats.
\newcommand{\qfig}[2][1]{\qf@check{#2}%
  \ifqf@full\qf@blk{#1}{#2}\else
  \ifnum\@nameuse{qf@wide@#2}=1 \expandafter\@firstoftwo\else\expandafter\@secondoftwo\fi
  {\end{multicols}\begin{center}\qf@img{#1}{\textwidth}{#2}\end{center}\begin{multicols}{2}}%
  {\qf@blk{#1}{#2}}\fi}
\newcommand{\qf@blk}[2]{\par\smallskip{\centering\qf@img{#1}{\linewidth}{#2}\par}\smallskip}
% \qfigi[scale]{id}: inline figure (inside option lists / tables / minipages); never breaks columns.
\newcommand{\qfigi}[2][1]{\qf@check{#2}\qf@img[valign=c]{#1}{\linewidth}{#2}}
\makeatother
% structure macros - lint reads these to check completeness and figure placement
\newcommand{\papersec}[1]{\par\bigskip{\centering\Large\bfseries PAPER-#1\par}\medskip}
\newcommand{\partsec}[2]{\par\medskip\noindent{\large\bfseries PART-#1: #2}\par\smallskip}
\newcommand{\sect}[1]{\par\smallskip\noindent\textbf{#1}\par}
\newcommand{\qn}[1]{\par\smallskip\noindent\textbf{#1.}~}
\newcommand{\paraq}[2]{\par\smallskip\noindent\textbf{Paragraph for Questions #1 and #2}\par}
\newenvironment{qoptions}{\begin{itemize}[leftmargin=1.8em,itemsep=1pt,topsep=2pt,label={}]}{\end{itemize}}
\newcommand{\opt}[2]{\item[(#1)] #2}
% figure registry (generated by score2_tools.py finalize)
\qfigdef{p1_phy_s1_q02_1}{Score-I - Test 18 (24 Apr 2022)-img/p1_phy_s1_q02_1.png}{83.7}{0}
\qfigdef{p1_phy_s1_q05_1}{Score-I - Test 18 (24 Apr 2022)-img/p1_phy_s1_q05_1.png}{179.5}{0}
\qfigdef{p1_phy_s2_q11_1}{Score-I - Test 18 (24 Apr 2022)-img/p1_phy_s2_q11_1.png}{72.2}{0}
\qfigdef{p1_phy_s3_q02_1}{Score-I - Test 18 (24 Apr 2022)-img/p1_phy_s3_q02_1.png}{98.0}{0}
\qfigdef{p1_phy_s3_q04_1}{Score-I - Test 18 (24 Apr 2022)-img/p1_phy_s3_q04_1.png}{335.1}{1}
\qfigdef{p1_chem_s1_q04_A}{Score-I - Test 18 (24 Apr 2022)-img/p1_chem_s1_q04_A.png}{33.7}{0}
\qfigdef{p1_chem_s2_q08_1}{Score-I - Test 18 (24 Apr 2022)-img/p1_chem_s2_q08_1.png}{152.1}{0}
\qfigdef{p1_chem_s2_q08_B}{Score-I - Test 18 (24 Apr 2022)-img/p1_chem_s2_q08_B.png}{64.4}{0}
\qfigdef{p1_chem_s2_q11_A}{Score-I - Test 18 (24 Apr 2022)-img/p1_chem_s2_q11_A.png}{59.0}{0}
\qfigdef{p1_chem_s2_q11_B}{Score-I - Test 18 (24 Apr 2022)-img/p1_chem_s2_q11_B.png}{50.8}{0}
\qfigdef{p1_chem_s2_q11_C}{Score-I - Test 18 (24 Apr 2022)-img/p1_chem_s2_q11_C.png}{32.3}{0}
\qfigdef{p1_chem_s3_q04_2}{Score-I - Test 18 (24 Apr 2022)-img/p1_chem_s3_q04_2.png}{168.2}{0}
\qfigdef{p1_chem_s3_q04_1}{Score-I - Test 18 (24 Apr 2022)-img/p1_chem_s3_q04_1.png}{153.3}{0}
\qfigdef{p2_phy_s1_q01_1}{Score-I - Test 18 (24 Apr 2022)-img/p2_phy_s1_q01_1.png}{122.5}{0}
\qfigdef{p2_phy_s2_q01_1}{Score-I - Test 18 (24 Apr 2022)-img/p2_phy_s2_q01_1.png}{70.6}{0}
\qfigdef{p2_phy_s2_q02_1}{Score-I - Test 18 (24 Apr 2022)-img/p2_phy_s2_q02_1.png}{160.2}{0}
\qfigdef{p2_phy_s2_q04_1}{Score-I - Test 18 (24 Apr 2022)-img/p2_phy_s2_q04_1.png}{196.8}{0}
\qfigdef{p2_phy_s2_q06_1}{Score-I - Test 18 (24 Apr 2022)-img/p2_phy_s2_q06_1.png}{131.0}{0}
\qfigdef{p2_phy_s3_q01_1}{Score-I - Test 18 (24 Apr 2022)-img/p2_phy_s3_q01_1.png}{103.2}{0}
\qfigdef{p2_phy_s3_q03_1}{Score-I - Test 18 (24 Apr 2022)-img/p2_phy_s3_q03_1.png}{118.8}{0}
\qfigdef{p2_phy_s3_q04_1}{Score-I - Test 18 (24 Apr 2022)-img/p2_phy_s3_q04_1.png}{148.0}{0}
\qfigdef{p2_phy_s3_q06_1}{Score-I - Test 18 (24 Apr 2022)-img/p2_phy_s3_q06_1.png}{67.6}{0}
\qfigdef{p2_chem_s2_q02_1}{Score-I - Test 18 (24 Apr 2022)-img/p2_chem_s2_q02_1.png}{144.0}{0}
\qfigdef{p2_chem_s2_q04_1}{Score-I - Test 18 (24 Apr 2022)-img/p2_chem_s2_q04_1.png}{220.9}{0}
\qfigdef{p2_chem_s2_q05_1}{Score-I - Test 18 (24 Apr 2022)-img/p2_chem_s2_q05_1.png}{316.5}{1}
\qfigdef{p2_chem_s3_q02_1}{Score-I - Test 18 (24 Apr 2022)-img/p2_chem_s3_q02_1.png}{177.5}{0}
\qfigdef{p2_chem_s3_q05_1}{Score-I - Test 18 (24 Apr 2022)-img/p2_chem_s3_q05_1.png}{206.9}{0}
% ---- body skeleton ----
% \begin{document}
% \begin{center}{\Large\bfseries JEE (Main+Advanced) Score-II --- Test NN (DD Mon YYYY)}\end{center}
% \begin{multicols}{2}
%   \papersec{1}   (only for two-paper tests)
%   \partsec{1}{PHYSICS} \sect{Section-I (i): Single Correct} \qn{1} ... \qfig{p1_phy_q01_1} ...
% \end{multicols}
% \end{document}

\begin{document}
\begin{center}{\Large\bfseries JEE (Main+Advanced) Score-I --- Test 18 (24 Apr 2022)}\end{center}
\raggedcolumns
\setlength{\tabcolsep}{3pt}
\begin{multicols}{2}
\papersec{1}
\partsec{1}{PHYSICS}
\sect{SECTION-I(i)}
\qn{1} Smoke consist of ash particles which acquire terminal velocity under influence of buoyant force and gravity very quickly after coming out of flame. If the steam engine of a train moving with constant velocity $\vec v_0$ releases smoke when wind is stationary :
\begin{qoptions}
\opt{A}{Smoke trail would be vertical and smoke particles would be travelling vertically as seen from ground.}
\opt{B}{Smoke trail would be at an acute angle to horizontal and smoke particles would be travelling along the smoke trail.}
\opt{C}{Smoke trail would be vertical but smoke particles would be travelling at an acute angle to the horizontal.}
\opt{D}{Smoke trail would be at an acute angle to horizontal but smoke particles would be travelling vertically.}
\end{qoptions}

\qn{2} A plane trapezium is in a plane containing principal axis of a lens. It's real image is formed. The image is of the shape of a square. The side length of square will be :-
\qfig{p1_phy_s1_q02_1}
\begin{qoptions}
\opt{A}{1.5 cm}
\opt{B}{2 cm}
\opt{C}{3 cm}
\opt{D}{6 cm}
\end{qoptions}

\qn{3} A parachutist descends uniformly at a speed of 8 m/s in calm weather. What will the terminal speed of the parachutist be if there is a 6 m/s crosswind in the horizontal direction?
\begin{qoptions}
\opt{A}{8 m/s}
\opt{B}{10 m/s}
\opt{C}{2 m/s}
\opt{D}{$2\sqrt7$ m/s}
\end{qoptions}

\qn{4} A particle of charge $+q$ and mass $m$ is in a uniform magnetic field $\vec B=B_0\hat z$. The initial position of the particle is $x=0$, $y=0$, $z=0$; the magnitude of the momentum of the particle is $p_0$. The particle is so projected that it crosses the point $\left(x=\dfrac{p_0}{qB_0},\,y=0,\,z=0\right)$. What is the time (in sec) after which it reaches this point? Take $\dfrac mq=\dfrac{12}{\pi}$ SI units and $B=0.2$ T.
\begin{qoptions}
\opt{A}{10}
\opt{B}{20}
\opt{C}{30}
\opt{D}{15}
\end{qoptions}

\qn{5} The impedance $Z_1$ in figure can be regarded as a pure resistance $R_1=12\ \Omega$, whereas the impedance $Z_2$ is associated with a series resistance $R_2=8\Omega$ and a capacitance $C=1\,\mu$F. If $f=5000$ Hz and $V_0=30$V, what is the power dissipated in $Z_2$ (in W)? Take $\pi^2=10$.
\qfig{p1_phy_s1_q05_1}
\begin{qoptions}
\opt{A}{2.57}
\opt{B}{3.12}
\opt{C}{257}
\opt{D}{None of these}
\end{qoptions}

\qn{6} Four students Ram, Shyam, Ghanshyam and Radheshyam perform an optical bench experiment to find focal length of a convex lens. Ram forms image of a distant luminous object on a small screen mounted on the optical bench. Other 3 students perform usual optical bench experiment. Their position of object needle, lens and image needle are given below. Which of the student is most accurate in his result?
\begin{center}\small
\begin{tabular}{|p{1.9cm}|p{1.4cm}|p{1.4cm}|p{1.4cm}|}
\hline
\centering Name & \centering Object & \centering Lens & \centering Image \tabularnewline
\hline
Ram & \_ & 15.0 cm & 30.0 cm\\
\hline
Shyam & 15.0 cm & 75.0 cm & 95.0 cm\\
\hline
Ghanshyam & 15.0 cm & 45.0 cm & 75.0 cm\\
\hline
Radheshyam & 15.0 cm & 60.0 cm & 82.5 cm\\
\hline
\end{tabular}
\end{center}
\begin{qoptions}
\opt{A}{Ram}
\opt{B}{Ghanshyam}
\opt{C}{Radheshyam}
\opt{D}{All are equally accurate}
\end{qoptions}

\sect{SECTION-I(ii)}
\qn{7} The de-broglie wavelength of 2 elementary particles in centre of mass frame is same and is equal to 5/8 times de-broglie wavelength of the lighter particle in lab frame and 5/2 times the de-broglie wavelength of the heavier particle in lab frame. Both particles have either parallel or anti parallel velocities. The possible combinations of particles are :
\begin{qoptions}
\opt{A}{$^{13}_{\,6}$C and $^{12}_{\,6}$C}
\opt{B}{$^{13}_{\,6}$C and $^{4}_{2}$He}
\opt{C}{$^{1}_{1}$H and $^{4}_{2}$He}
\opt{D}{$^{7}_{3}$Li and $^{13}_{\,6}$C}
\end{qoptions}

\qn{8} A car has side window made of glass (shear strength 40 M Pa) having size 40 cm $\times$ 40 cm. The car is travelling with a velocity of 40 m/s in still air. Density of air can be assumed to be 1.25 kg/m$^3$. The windows are closed.
\begin{qoptions}
\opt{A}{If the glass is very thin, it may break and fall into the car.}
\opt{B}{If the glass is very thin, it may break \& fall out of the car.}
\opt{C}{If we want that glass should not break, it's thickness should be greater than 2.5 $\mu$m.}
\opt{D}{If we want that glass should not break, it's thickness should be greater than 8.75 $\mu$m.}
\end{qoptions}

\qn{9} Two sound waves propagate in a medium. Their equation for displacement are given by

$S_1=1\mu\text{m}\sin\left(100\pi t-\dfrac\pi4x+\dfrac\pi3\right)$ ; $S_2=1\mu\text{m}\sin\left(\dfrac\pi4x-100\pi t\right)$

Which of the following statement is/are \textbf{CORRECT} ?
\begin{qoptions}
\opt{A}{The velocity of both waves is the same.}
\opt{B}{Both the waves are equally intense.}
\opt{C}{They superimpose to produce a wave of amplitude 1$\mu$m.}
\opt{D}{Both waves are plane waves, propagating in +ve x direction.}
\end{qoptions}

\qn{10} When heat $Q=600$ J is added to a mixture of helium and nitrogen at a constant volume the temperature of the mixture increases by 15 K and if the same amount of heat is given to the same mixture at constant pressure, the temperature of the mixture will increase by 10 K. Take $R=\dfrac{25}{3}$ Jmol k.
\begin{qoptions}
\opt{A}{Ratio of number of nitrogen molecules and helium molecules in mixture is 1 : 1.}
\opt{B}{The internal energy changes more in isochoric process.}
\opt{C}{Work done by the gas is more in isobaric process.}
\opt{D}{The total number of moles in the mixture is 2.4 moles.}
\end{qoptions}

\qn{11} An electric dipole and point charge $-q$ are kept on $x$ axis at a distance $a$ from each other. Assume $a$ is much larger than size of dipole. It is seen that $\vec E$ vanishes at a height $y$ above $-q$.
\qfig{p1_phy_s2_q11_1}
\begin{qoptions}
\opt{A}{$y=\sqrt2\,a$}
\opt{B}{$p=\dfrac{3\sqrt3}{2\sqrt2}\,qa$}
\opt{C}{Potential energy of interaction of dipole with $-q$ is $\dfrac{-9\sqrt3\,k}{4a}q^2$}
\opt{D}{Given situation is not possible.}
\end{qoptions}

\qn{12} In a meter bridge experiment, when we measure the resistance of two resistors connected in series, null point is obtained at 80 cm. When they are in parallel, null point is at 40 cm. What are the approximate null points when each of the resistances are measured ? Assume that resistance of resistance box remains the same.
\begin{qoptions}
\opt{A}{76 cm}
\opt{B}{60 cm}
\opt{C}{46 cm}
\opt{D}{28 cm}
\end{qoptions}

\sect{SECTION-II}
\qn{1} In a sample of $^{211}_{\,86}$Rn nuclei, 25 percent decay by $\alpha$-decay to $^{207}_{\,84}$Po and the rest decay by electron capture to $^{211}_{\,85}$At. What is the ratio of the mean life for $\alpha$-decay, $\tau_\alpha$, to the mean life for electron capture, $\tau_\beta$ ?

\qn{2} Two simple pendulum made of small masses $m$ each are pulled to side as shown and released. The collision is elastic. After what time (in sec) from the release will the left mass first reach point from where it was released ? Length of both strings is 1 m, take $g=\pi^2$ m/s$^2$.
\qfig{p1_phy_s3_q02_1}

\qn{3} Two wavelengths $\lambda_1$ and $\lambda_2$ are used in the double-slit experiment. If $\lambda_1$ is 430 nm, what value must the other have for the fourth-order bright fringe of $\lambda_1$ to fall on the sixth-order bright fringe of the $\lambda_2$ ?

\qfullbegin
\qn{4} As seen from figure, currents $I_1$ and $I_2$ can be independently increased, decreased, or held constant by varying resistances $R_1$ and $R_2$. Let $L_1=50$ mH, $L_2=40$ mH, $M=15$ mH. $I_1$ is increased at the rate of 120 A/s and $I_2$ is decreased at the rate of 200 A/s. Compute $|v_1|$ in volt. (Induced emf in coil-1)
\qfig{p1_phy_s3_q04_1}
\qfullend

\qn{5} A light ray is incident on one of the faces of a glass prism perpendicularly to it. The angle of emergence of the light ray on another face of the prism is $45^\circ$. What is the apex angle $A$ of the prism if the speed of light in this glass prism is 240000 km/s? Fill $\sin A$ in OMR sheet.

\qn{6} Two conducting plates, each 3 cm $\times$ 6 cm, and three slabs of dielectric, each 1 cm $\times$ 3 cm $\times$ 6 cm and having dielectric constants of 1, 2, and 3, are assembled into a parallel plate capacitor with $d=3$ cm. Determine the two values of capacitance obtained by the two possible methods of assembling the capacitor. Find $\dfrac{(C_{\text{higher}}-C_{\text{Lower}})}{\epsilon_0}$ (in cm).

\partsec{2}{CHEMISTRY}
\sect{SECTION-I(i)}
\qn{1} The \textbf{CORRECT} order of first ionization potential among the following is :-
\begin{qoptions}
\opt{A}{F$^-$ $>$ Cl$^-$}
\opt{B}{Ga $>$ T$\ell$}
\opt{C}{Sn $<$ Pb}
\opt{D}{O$^-$ $>$ S$^-$}
\end{qoptions}

\qn{2} At isoelectric point, the amino acid has-
\begin{qoptions}
\opt{A}{Least viscosity}
\opt{B}{Maximum surface tension}
\opt{C}{Maximum solubility}
\opt{D}{All of these}
\end{qoptions}

\qn{3} Sodium salts are highly soluble in water and hence undergoes complete ionisation in water. If the freezing point of 0.35 molal aq. Na$_x$A solution is $-$ 3.255\textdegree C, the value of `x' is ($K_f$ = 1.86 K--kg/mole)
\begin{qoptions}
\opt{A}{2}
\opt{B}{3}
\opt{C}{4}
\opt{D}{5}
\end{qoptions}

\qn{4} The specie among the following which when brought near a charged comb, will attract it.
\begin{qoptions}
\opt{A}{\qfigi{p1_chem_s1_q04_A}}
\opt{B}{IF$_7$}
\opt{C}{SO$_3$}
\opt{D}{C$_6$H$_6$}
\end{qoptions}

\qn{5} m-chlorobenzaldehyde on reaction with conc. KOH at room temperature gives $-$
\begin{qoptions}
\opt{A}{Potassium m-chlorobenzoate and m hydroxybenzaldehyde}
\opt{B}{m-hydroxy benzaldehyde and m-chlorobenzyl alcohol}
\opt{C}{m-chlorobenzyl and m-hydroxybenzyl alcohol.}
\opt{D}{Potassium m-chlorobenzoate and m-chlorobenzyl alcohol.}
\end{qoptions}

\qn{6} Calculate the minimum volume of 1M-HNO$_3$ solution required to completely dissolve 95.25 gm Cu, producing NO. (Cu = 63.5)
\begin{qoptions}
\opt{A}{4.00 L}
\opt{B}{1.00 L}
\opt{C}{3.00 L}
\opt{D}{3.60 L}
\end{qoptions}

\sect{SECTION-I(ii)}
\qn{7} Which of the following statement(s) regarding REFINING is/are \textbf{CORRECT} ?
\begin{qoptions}
\opt{A}{Zone refining is based on the principle that impurities are more soluble in molten state}
\opt{B}{In liquation, the M.P. of metal should be less than that of impurities}
\opt{C}{In electrolytic refining, the more basic metal remains in the solution and the less basic ones go to anode mud.}
\opt{D}{In Van-Arkel method, the metal iodide is decomposed on a Tungsten filament, depositing pure metal at the filament}
\end{qoptions}

\qn{8} Which of the following are linked to the reaction ?
\qfig{p1_chem_s2_q08_1}
\begin{qoptions}
\opt{A}{carbocation rearrangement}
\opt{B}{\qfigi{p1_chem_s2_q08_B}}
\opt{C}{Friedel-Crafts reaction}
\opt{D}{Methyl shift}
\end{qoptions}

\qn{9} A lead storage cell is discharged which causes the H$_2$SO$_4$ electrolyte to change from a concentration of 40\% by weight (density = 1.260 gm/ml) to 28\%, by weight. The original volume of electrolyte was one litre. Identify the correct statement(s) :-
\begin{qoptions}
\opt{A}{The overall cell reaction is: Pb(s) + PbO$_2$(s) + 2H$_2$SO$_4$(aq) $\to$ 2PbSO$_4$(s) + 2H$_2$O($\ell$).}
\opt{B}{A total of 2.0 moles of H$_2$SO$_4$ have reacted.}
\opt{C}{The total charge released from anode of the cell is $1.93\times10^5$ coulomb.}
\opt{D}{The mass of electrolytic solution has decreased.}
\end{qoptions}

\qn{10} Which of the following reactions can be used to produce Boric Anhydride ?
\begin{qoptions}
\opt{A}{B$_2$H$_6$ + 2CO $\longrightarrow$}
\opt{B}{H$_3$BO$_3$ $\xrightarrow{\Delta,\ <\ 370\text{K}}$}
\opt{C}{Na$_2$[B$_4$O$_5$(OH)$_4$].8H$_2$O $\xrightarrow[\text{heating}]{\text{strong}}$}
\opt{D}{B$_2$H$_6$ + 3O$_2$ $\longrightarrow$}
\end{qoptions}

\qn{11} Which of the following compound can not exhibit optical isomerism ?
\begin{qoptions}
\opt{A}{\qfigi{p1_chem_s2_q11_A}}
\opt{B}{\qfigi{p1_chem_s2_q11_B}}
\opt{C}{\qfigi{p1_chem_s2_q11_C}}
\opt{D}{CH$_3$--CH=C=CH$_2$}
\end{qoptions}

\qn{12} For H-atom, the wave function for an electron is $\psi=\dfrac{1}{81}\sqrt{\dfrac{2}{\pi a_0^2}}\cdot\sigma\cdot(6-\sigma)\cdot e^{-\sigma/2}\cdot\cos\theta$, where $\sigma=\dfrac{2Zr}{na_0}$ and $a_0=0.529$ \AA. Select the \textbf{CORRECT} information(s) about the orbital to which this electron belongs.
\begin{qoptions}
\opt{A}{3p$_z$}
\opt{B}{Only one angular node}
\opt{C}{Only one radial node}
\opt{D}{Nodal plane is XY plane}
\end{qoptions}

\sect{SECTION-II}
\qn{1} Number of reactions that will produce \textbf{BLACK} precipitate among the following.\\
(i) 2Ag$^+$ + H$_2$S $\longrightarrow$\\
(ii) Na$_2$SnO$_2$ + Bi$^{+3}$ $\xrightarrow{\text{OH}^-}$\\
(iii) 2Bi$^{3+}$ + 3H$_2$S $\longrightarrow$\\
(iv) Bi$^{3+}$ + 3OH$^-$ $\longrightarrow$\\
(v) [Cd(CN)$_4$]$^{2-}$ + H$_2$S $\longrightarrow$\\
(vi) AsO$_4^{3-}$ + 12MoO$_4^{2-}$ + 3NH$_4^+$ + 24H$^+$ $\to$\\
(vii) Fe$^{2+}$ + 2K$^+$ + [Fe(CN)$_6$]$^{4-}$ $\longrightarrow$

\qn{2} The total number of alkenes possible by dehydrobromination of 3-bromo-3-cyclopentylhexane using alcoholic KOH (by E$_2$ mechanism) is.

\qn{3} The rate of the reaction gets doubled when the temperature changes from 7\textdegree C to 17\textdegree C. By what factor (times) will it change for the temperature change from 17\textdegree C to 27\textdegree C ?

[Antilog (0.30) = 2.0, Antilog (0.280) = 1.91, Antilog (0.322) = 2.1]

\qn{4} Number of splitting configuration(s) against the given complex that is/are \textbf{INCORRECTLY} indicated.
\qfig{p1_chem_s3_q04_2}
\qfig{p1_chem_s3_q04_1}

\qn{5} Number of moles of CrO$_3$ required to oxidise 6 moles of 2-heptanol :

\qn{6} Calculate $\Delta G$ (in kJ/mol) for the change :

H$_2$(g, 300K, 1 bar) $\to$ H$_2$(g, 600K, 2 bar)

Given that H$_2$(g) is ideal and for H$_2$(g), $C_{p,m}$ = 28.314 J/K-mol, $S_{300\text{K}}$, 1 bar = 130.0 J/K-mol. ($\ell$n2 = 0.70)

\partsec{3}{MATHEMATICS}
\sect{SECTION-I(i)}
\qn{1} Let $a_1,a_2,a_3,a_4\ldots a_n,a_{n+1},a_{n+2}$ forms an arithmetic progression. If $S_n=\displaystyle\sum_{i=1}^{n}\dfrac{a_ia_{i+1}a_{i+2}}{a_i+a_{i+2}}$ and $a_1=5$ and common difference of progression is 3 then value of $S_{10}$ is :-
\begin{qoptions}
\opt{A}{2637}
\opt{B}{2637.50}
\opt{C}{2636}
\opt{D}{2638}
\end{qoptions}

\qn{2} A variable chord PQ of parabola $y^2=4x$ is drawn parallel to a line $y=x$. Then which of the following point lies on the locus of point of intersection of the normals at P and Q.
\begin{qoptions}
\opt{A}{(7, 2)}
\opt{B}{(7, 1)}
\opt{C}{(3, 6)}
\opt{D}{(10, 2)}
\end{qoptions}

\qn{3} If $\displaystyle\lim_{x\to0^+}\left(\dfrac{\sin^2x}{\left[\frac x\pi\right]+\frac{x^2}{\pi^2}}\right)\times\dfrac{\sin(\sin x)-\sin x}{ax^5+bx^3+c}=\dfrac{-\pi^2}{12}$. Then value of $b+c$ is

(Where [.] denotes greatest integer function)
\begin{qoptions}
\opt{A}{2}
\opt{B}{3}
\opt{C}{4}
\opt{D}{$-2$}
\end{qoptions}

\qn{4} The coefficient of $x^{70}$ in the product $(x-1)(x^2-2)(x^3-3)\ldots(x^{12}-12)$ is :
\begin{qoptions}
\opt{A}{4}
\opt{B}{5}
\opt{C}{6}
\opt{D}{7}
\end{qoptions}

\qn{5} If $\hat i\times\left[(\vec a-\hat j)\times\hat i\right]+\hat j\times\left[(\vec a-\hat k)\times\hat j\right]+\hat k\times\left[(\vec a-\hat i)\times\hat k\right]=0$ and let $\vec b=2\left(\hat i+\hat j+\hat k\right)$ then $\left(\vec a\cdot\vec b\right)$ is
\begin{qoptions}
\opt{A}{1}
\opt{B}{2}
\opt{C}{3}
\opt{D}{4}
\end{qoptions}

\qn{6} If $z\in C$ then area enclosed by the curve $2\left|z^2+|z|^2\right|+3\left|z^2-|z|^2\right|=6|z|$, $z\neq0$ is (Where $C$ denotes set of complex number)
\begin{qoptions}
\opt{A}{3 unit$^2$}
\opt{B}{6 unit$^2$}
\opt{C}{$3\pi$ unit$^2$}
\opt{D}{$6\pi$ unit$^2$}
\end{qoptions}

\sect{SECTION-I(ii)}
\qn{7} If area enclosed by the graph of function $y=\ell n(x+1)$, $y$-axis and the line $y=1$ is $A$, then choose correct option(s)- (where [.] and \{.\} denotes greatest integer function and fractional part function respectively)
\begin{qoptions}
\opt{A}{\{A\} = e $-$ 2}
\opt{B}{[A] = 1}
\opt{C}{[A] = 0}
\opt{D}{[$-$A] = $-$1}
\end{qoptions}

\qn{8} If a variable tangent of the circle $x^2+y^2=1$ intersect the ellipse $x^2+2y^2=4$ at the points P and Q. If locus of point of intersection of tangents at P and Q is a conic C. Then which of the following is are true
\begin{qoptions}
\opt{A}{C is hyperbola}
\opt{B}{C is ellipse}
\opt{C}{eccentricity of C is $\dfrac{\sqrt3}{2}$}
\opt{D}{eccentricity of C is 1}
\end{qoptions}

\qn{9} ABC is a triangle in which $\angle C=2\angle B$. D is a point on BC such that AD bisects $\angle BAC$ and $AC=BD$. Then which of the following is/are true
\begin{qoptions}
\opt{A}{$\cos B+\cos C=\dfrac{\sqrt5}{2}$}
\opt{B}{$\cos B-\cos C=\dfrac12$}
\opt{C}{$\dfrac{\sin A}{\sin C}=1$}
\opt{D}{$\dfrac{\sin A}{\sin C}=\dfrac{\sqrt5+1}{4}$}
\end{qoptions}

\qn{10} In kota city a person owns independently a mercedes car with probability $\dfrac{3}{10}$ and Audi with probability $\dfrac{4}{10}$. If he has mercedes only then he keeps a driver with probability $\dfrac{6}{10}$ where as if he owns Audi car only, then he keeps a driver with probability $\dfrac{7}{10}$ where as if he keeps both type of cars then his probability of keeping a driver is $\dfrac{9}{10}$. Then
\begin{qoptions}
\opt{A}{Probability that person keeps a driver is $\dfrac{412}{1000}$}
\opt{B}{Probability that person keeps a driver is $\dfrac{71}{125}$}
\opt{C}{Given that person keeps driver probability that he owns Audi car is $\dfrac{54}{103}$}
\opt{D}{Given that person keeps driver, the probability that he own Audi car is $\dfrac{76}{103}$}
\end{qoptions}

\qn{11} If the projection of the line $\dfrac x2=\dfrac{y-1}{2}=\dfrac{z-1}{1}$ on a plane P is $\dfrac x1=\dfrac{y-1}{1}=\dfrac{z-1}{-1}$. Then
\begin{qoptions}
\opt{A}{Equation of plane P is $x+y+2z-3=0$}
\opt{B}{Equation of plane P is $x+y+2z+3=0$}
\opt{C}{Distance of plane P from origin is $\sqrt{\dfrac32}$}
\opt{D}{Distance of plane P from point $(1,1,-1)$ is $\dfrac{3}{\sqrt6}$}
\end{qoptions}

\qn{12} If Rolle's theorem is applicable to the function $f$ defined by $f(x)=\begin{cases}ax^2+b & ;\ |x|<1\\ 1 & ;\ |x|=1\\ \dfrac{c}{|x|} & ;\ |x|>1\end{cases}$ in the interval $[-3,3]$ then which of the following alternative is/are \textbf{CORRECT}.
\begin{qoptions}
\opt{A}{$a+b+c=2$}
\opt{B}{$|a|+|b|+|c|=3$}
\opt{C}{$2a+4b+3c=8$}
\opt{D}{$4a^2+4b^2+4c^2=14$}
\end{qoptions}

\sect{SECTION-II}
\qn{1} If $A$ and $B$ are two matrices of order 3 such that $2A+3BB^T=I$ and $B^{-1}=A^T$. Let $\alpha=\det(A^{-1}-3B^3+BA)$ and $\beta=\text{Tr}(A^{-1}-3B^3-BA)$, then value of $\dfrac\alpha\beta$ is :

\qn{2} $f(x)$ is polynomial which satisfies $|f(x)|\le|e^x-1|$ $\forall\ x\in\mathbb R$ and $A=\{x:x$ is all possible values of $f'(0)\}$, then number of integer in set $A$ is :

\qn{3} Given two circle $x^2+y^2+3\sqrt2\,(x+y)=0$ and $x^2+y^2+5\sqrt2\,(x+y)=0$, the radius of third circle which is tangent to the two given circles and to their common diametre is :-

\qn{4} If $N$ is the number of ways in which 5 distinct balls can be distributed among 3 people such that no person is left empty handed, then $\dfrac{N}{10}$ is equal to

\qn{5} Let $F(x)$ be a non-negative continuous function defined on $R$ such that $F(x)+F\left(x+\dfrac12\right)=1$. Then value of $\dfrac{1}{1000}\displaystyle\int_0^{1500}F(x)\,dx$.

\qn{6} If $(x,y)$ satisfy the equation $\tan^4x+\tan^4y+2\cot^2x\cot^2y=3+\sin^2(x+y)$ sum of all the possible values of $x$ where $x,y\in(0,2\pi)$
\papersec{2}
\partsec{1}{PHYSICS}
\sect{SECTION-I}
\qn{1} An infinite current carrying wire is placed at a distance $b$ from a very small square loop of side $a$. The current in loop is $i$. Neglect any induced emf and assume these currents remain constant.
\qfig{p2_phy_s1_q01_1}
\begin{qoptions}
\opt{A}{The loop is attracted to the wire by a force $F=\dfrac{\mu_0Ii}{2\pi b^2}a^2$.}
\opt{B}{The loop must be projected away from wire with a kinetic energy of $\dfrac{\mu_0Iia^2}{3\pi b}$, so that it does not return back.}
\opt{C}{If the loop is projected into the plane with a kinetic energy of $\dfrac{\mu_0Iia^2}{4\pi b}$, it moves in a circular path.}
\opt{D}{If the loop is projected into the plane with a kinetic energy of $\dfrac{\mu_0Iia^2}{2\pi b}$, it moves in a circular path.}
\end{qoptions}

\qn{2} A simple pendulum time period at night 12'O clock is 2 sec if it is on equator. In which of the following cases will it's time period be greater than 2 sec?
\begin{qoptions}
\opt{A}{It is taken to poles.}
\opt{B}{The temperature is increased.}
\opt{C}{It is taken on a geostationary satellite.}
\opt{D}{It is taken to a high hill on equator.}
\end{qoptions}

\qn{3} A, B and C are the mean positions of three particle in a line along which Simple harmonic waves of amplitude 10 units is travelling. At a certain instant particle A has a displacement of 8.66 decreasing, B has displacement of $+5$ increasing, and C has displacement of $+8.66$ decreasing. The distance AB = 13 cm and is less than one wavelength. BC is also less than one wavelength. Assume $y=A\sin(kx-\omega t)$.
\begin{qoptions}
\opt{A}{The phase angles of A is $2\pi/3$ and that of B is $13\,\pi/6$.}
\opt{B}{The wavelength is 52 cm.}
\opt{C}{The distance BC is 39 cm.}
\opt{D}{The position of particles D and E in between them, for each of which the displacement is zero is $29/9$ cm, $169/18$ cm from A.}
\end{qoptions}

\qn{4} The rate at which energy is produced in mammals is proportional to their volume. Assume that mammals lose energy according to Newton's Law of cooling. Their body temperature are 27\textdegree C and surrounding is at 7\textdegree C. Assume that each mammal is curled up to form a sphere of different radius.

(i) Mouse of radius 2 cm.

(ii) human being of radius 40 cm.

(iii) elephant of radius 2 m.

The energy intake/mass to maintain their temperature are $E_1$, $E_2$ \& $E_3$ respectively. Assume same emissivity and density.
\begin{qoptions}
\opt{A}{$E_1>E_2>E_3$}
\opt{B}{$E_3>E_2>E_1$}
\opt{C}{rate of heat loss/area for all of them is same.}
\opt{D}{rate of heat loss/area is largest for elephant.}
\end{qoptions}

\qn{5} A car of mass 1200 kg started from rest and speeded up at an acceleration of 2 m/s$^2$ along a straight horizontal path of length 200 m. Its wheels did not slide. (The mass of the wheels can be neglected.)
\begin{qoptions}
\opt{A}{The total frictional force exerted between the ground and the wheels was 2400 N.}
\opt{B}{The final kinetic energy of the car is 480 kJ.}
\opt{C}{Work done by the frictional force is 240 kJ.}
\opt{D}{Final speed is 72 km/hr.}
\end{qoptions}

\qn{6} A shiny spoon is held in front of our eye at a distance of 25 cm such that the stem of the spoon is vertical. If the concave side of the spoon is observed, then an inverted image of our head can be seen, whilst if the convex part of the spoon is observed, then the image is upright. The radius of curvature of the vertical section of the spoon is 5 cm.
\begin{qoptions}
\opt{A}{The upright image is shorter than the inverted image.}
\opt{B}{The inverted image is shorter than the upright image.}
\opt{C}{The inverted image is at a distance of $\dfrac{200}{9}$ cm from us.}
\opt{D}{The upright image is at a distance of $\dfrac{250}{11}$ cm from us.}
\end{qoptions}

\sect{SECTION-II}
\qn{1} Consider a flat plate of area $S$ moving horizontally through a medium having stationary particles of mass $m$ each. The number of particles/volume in medium is $n$. The collision of particles with plate is elastic. If angle $\alpha$ is varied, the maximum force on the plate in the vertical direction is $knmu^2S$. Find the value of $k$.
\qfig{p2_phy_s2_q01_1}

\qn{2} The picture shown below a picture of vernier screw gauge. Main scale or pitch is 0.5 mm. There are 50 divisions on circular scale. 9 circular scale divisions correspond to 10 vernier scale divisions. What is the reading $R$ (in mm)? Fill $10R$ in the OMR.
\qfig{p2_phy_s2_q02_1}

\qn{3} An optical system contains a thin lens, $n=1.4$, with positive curvature of radius $R_1=25$ cm for surface 1 and negative curvature of radius $R_2=-25$ cm for surface 2. This lens collects light from the right side. To the left of the lens we want to keep a flat plate of thickness $t$ of the same glass so that the light that radiates from a distant object is focused on a screen 35 cm to the left of the lens. What is thickness $t$ (in cm) ?

\qn{4} A drop of liquid of radius $r$, density $\rho$ and surface tension $\sigma$ oscillates in gravity free space as shown in figure. If the time period of oscillation of this drop is 2 sec, what is the time period (in sec) of oscillation of a drop of radius $2r$. The second drop is made of same liquid as first.
\qfig{p2_phy_s2_q04_1}

\qn{5} Suppose that two electrons are in orbit around one proton (an H ion), both in an $n=1$ level. Assuming standard reference for potential energy, find the modulus of mechanical energy of this system (in eV). Take energy of ionization of Bohr atom as 13.6 eV. Assume that the total angular momentum is quantized as in Bohr's theory.

\qn{6} On perfectly slippery ice, two identical boards lie in contact (figure). At the left edge a rough block is placed on the first board. When it was pushed, it reached the right edge of the second board and stayed on it. The final speed of second board is $n$ times the final speed of first board. What is the value of $n$? Weight and size of block is much less than mass and the length of the boards.
\qfig{p2_phy_s2_q06_1}

\sect{SECTION-III}
\qn{1} Half portion of a spherical capacitor is filled with a dielectric of dielectric constant $k=2$ and conductivity $\sigma$. The charge given to spherical capacitor is $Q_0$. Due to the conductivity of dielectric charge leaks and the time constant for the discharge circuit is $x$ time of $\dfrac{\epsilon_0}{\sigma}$. Find the value of $x$.
\qfig{p2_phy_s3_q01_1}

\qn{2} At low temperatures, the molar heat capacity of solids at constant pressure depends on the absolute temperature T according to the law $C_T=C_{T_0}\left(\dfrac{T}{T_0}\right)^3$. One mole (40 g) of solid argon at temperature 8K is brought into thermal contact with two kilograms of solid argon at a temperature of 1K and they insulated everything together. What will be the temperature (in K) of the substance when thermal equilibrium is established? Round off to nearest integer.

\qn{3} A body of mass 1 kg is suspended by means of a spring and a thread between two walls as shown in the figure. The measure of both angles between the thread and the horizontal and the spring and the horizontal is shown in the figure. What is the initial acceleration (in m/s$^2$) of the object if the thread breaks ?
\qfig{p2_phy_s3_q03_1}

\qn{4} In a photoelectric experimental setup ultraviolet light of wavelength 350 nm is incident on emitter plate (E). The work function of the emitter plate is $\phi=2.2$ eV. AB is a uniform wire resistor having length $L=100$ cm. Resistance of the wire AB is 10 $\Omega$ and the emf of the battery is $V=10$ volt. The sliding contact J can be moved along the wire AB. When the sliding contact is placed at end B of the wire resistor the micro-ammeter shows a reading of $i=6\ \mu$A. Assume that photoelectric current is very small compared to the current through the cell. Find the reading of the ammeter (in $\mu$A) when the slider is moved to end A of the wire.
\qfig{p2_phy_s3_q04_1}

\qn{5} In the manufacture of boiler filament we used a Nichrome wire of volume 10 cm$^3$. What is the mass of water (in multiples of 10 gm/min) that can be heated every minute from the temperature of 10\textdegree C to a temperature of 100\textdegree C. The current density in it = 3 A/mm$^2$. Boiler efficiency = 70\%. Resistivity of nichrome is 1 $\mu\Omega$m. $S_{\text{Water}}$ = 4200 J/kg\textdegree C.

\qn{6} We have a fixed smooth insulating ring of radius $R$ kept in a horizontal plane. A bead of mass $m$ and charge $+q$ is threaded on the ring. A time varying magnetic field $B=B_0+\alpha t$ is present. It is perpendicular to the plane of ring as shown. Find the force exerted (in $10^{-4}$N) on the bead by the ring at $t=5$ sec. Take $B_0=0.01$ T, $\alpha=10^{-3}$ T/sec., $R=1$ m, $q=1$ mC, $m=\dfrac{10^{-6}}{16}$ kg.
\qfig{p2_phy_s3_q06_1}

\partsec{2}{CHEMISTRY}
\sect{SECTION-I}
\qn{1} Consider the given electronic configuration of atom/ion :-

$1s^0\ 2s^0\ 2p^0\ 3s^0\ 3p^0\ 3d^0\ 4s^1$

Select \textbf{INCORRECT} statement(s) :-
\begin{qoptions}
\opt{A}{Electron is in 3$^{\text{rd}}$ excited state of given atom/ion}
\opt{B}{Degeneracy of energy level in which electron is present in given configuration = 32}
\opt{C}{Total number of quantum number sets without changing energy of electron = 16}
\opt{D}{Electron is present in an non directional orbital}
\end{qoptions}

\qn{2} p-Chloroaniline and anilinium hydrochloride can be distinguished by :-
\begin{qoptions}
\opt{A}{Sandmeyer reaction}
\opt{B}{NaHCO$_3$}
\opt{C}{AgNO$_3$}
\opt{D}{Carbylamine test}
\end{qoptions}

\qn{3} Select the correct statement(s):
\begin{qoptions}
\opt{A}{Heating is the transfer of energy that makes use of disorderly, apparently random, molecular motion in the surroundings.}
\opt{B}{Work is the transfer of energy that makes use of ogranised motion of the molecules}
\opt{C}{The internal energy of an isolated system is constant.}
\opt{D}{A reversible change in thermodynamics is a change that can be reversed by an infinitesimal modification of a variable.}
\end{qoptions}

\qn{4} For an element X$_{(g)}$ value of negative electron gain enthalpy is second highest in periodic table elements. Select \textbf{CORRECT} statement(s) for element X.
\begin{qoptions}
\opt{A}{Electronegativity of X is higher as compared to Cl$_{(g)}$ on Paulings scale.}
\opt{B}{Negative electron gain enthalpy of X$^+_{(g)}$ is more than negative electron gain enthalpy of O$^+_{(g)}$.}
\opt{C}{Negative electron gain enthalpy of X$_{(g)}$ is more than negative electron gain enthalpy of Cl$_{(g)}$.}
\opt{D}{Ionisation energy of X$^-_{(g)}$ is less than ionisation of Cl$^-_{(g)}$}
\end{qoptions}

\qn{5} Which of the following reaction are used to confirm presence of carbonyl group?
\begin{qoptions}
\opt{A}{Reaction with hydroxylamine}
\opt{B}{Reaction with hydrazine Sol}
\opt{C}{Reaction with phenylhydrazine}
\opt{D}{Reaction with semicarbazide hydrochloride}
\end{qoptions}

\qn{6} Identify the statement(s) which is/are \textbf{NOT CORRECT} with respect to the surface phenomenon.
\begin{qoptions}
\opt{A}{If on adding electrolyte in an emulsion, the conductivity decreases, then it will be an oil in water type emulsion.}
\opt{B}{Tyndall effect is observed when the refractive indices of the dispersed phase and dispersion medium differ largely.}
\opt{C}{Macromolecular colloids are generally lyophobic in nature.}
\opt{D}{Gases which can react with the adsorbents generally show chemisorption.}
\end{qoptions}

\sect{SECTION-II}
\qn{1} XeF$_6$ is not stored in glass bottles due to the formation of P and Q compounds. Here P is a nonpolar molecule and Q is a polar molecule formed by reaction of two mole of XeF$_6$ and one mole of acidic oxide present in glass.

Find the value of $(y-x)$

Here $x$ = number of similar bond angles in P molecule

\hspace{1.2em}$y$ = Number of d-orbitals used for bonding by central atom of Q

\qn{2} \qfigi[0.85]{p2_chem_s2_q02_1}\par
Number of $\alpha$-Hydrogen atoms in Hydrocarbon formed as major product.

\qn{3} The enthalpy of formation of KCl(s) (in Kcal/mol) from the following data is\\
(i) KOH(aq) + HCl(aq) $\to$ KCl(aq) + H$_2$O($\ell$); $\Delta H=-13.7$ kcal\\
(ii) H$_2$(g) + $\tfrac12$O$_2$(g) $\to$ H$_2$O($\ell$); $\Delta H=-68.4$ kcal\\
(iii) $\tfrac12$H$_2$(g) + $\tfrac12$Cl$_2$(g) + aq $\to$ HCl(aq); $\Delta H=-39.3$ kcal\\
(iv) K(s) + $\tfrac12$O$_2$(g) + $\tfrac12$H$_2$(g) + aq $\to$ KOH(aq); $\Delta H=-116.5$ kcal\\
(v) KCl(s) + aq $\to$ KCl(aq); $\Delta H=+4.4$ kcal

\qn{4} \qfigi[0.8]{p2_chem_s2_q04_1}\par
Find the sum of sigma and pi bonds in anionic part of (R).

\qfullbegin
\qn{5} Identify total number compounds which are soluble in aqueous NaOH as well as NaHCO$_3$?
\qfig[0.85]{p2_chem_s2_q05_1}
\qfullend

\qn{6} A solution contains 0.1M $-$ Ba$^{2+}$ and 0.2 M $-$ Sr$^{2+}$. If solid Na$_2$CrO$_4$ is added in this solution, what is the maximum percentage of Ba$^{2+}$ precipitated as BaCrO$_4$, before any precipitation of SrCrO$_4$, is

[$K_{sp}$ (BaCrO$_4$) = $1.2\times10^{-10}$ ; $K_{sp}$ (SrCrO$_4$) = $4.0\times10^{-9}$]

\sect{SECTION-III}
\qn{1} Find out the number of sets, in which first specie would be expected to have the higher C--O stretching frequency?\\
(CO, CO$_2$),\\
([Ti(CO)$_6$]$^{2\ominus}$, [V(CO)$_6$]$^{\ominus}$)\\
(Cr(CO)$_6$, [Mn(CO)$_6$]$^+$)\\
([Co(CO)$_4$]$^{\ominus}$, [Ni(CO)$_4$])\\
([(PCl$_3$)$_3$Mo(CO)$_3$], [(PMe$_3$)$_3$Mo(CO)$_3$])

\qn{2} How many alkenes are produced when following compound is heated with alc. KOH :
\qfig{p2_chem_s3_q02_1}

\qn{3} The density of crystalline sodium chloride is 2.167 gm/cm$^3$. What would be the approximate side length (in cm) of cube containing one mole of NaCl ?

\qn{4} Consider the following species.

OCl$_2$, N(SiH$_3$)$_3$, CCl$_3^{\ominus}$, $:$CCl$_2$, CO, (GeH$_3$)NCS, B(OMe)$_3$

Find the product of \textbf{(x.y)} if

x = number of species in which $2p_\pi$--$3d_\pi$ back bond present

y = number of species in which direction of lone pair donation is from 2$^{\text{nd}}$ period element to 3$^{\text{rd}}$ period element

\qn{5} \qfigi{p2_chem_s3_q05_1}\par

\qn{6} There are two vessels of same volume consisting same number of moles of two different gases at the same temperature. One of the gas is CH$_4$ and the other is unknown gas, G. Assuming that the molecules of G are under random motion whereas in CH$_4$, except one, all are stationary. The collision diameter is same for both the gases and the average speed of CH$_4$ is $\sqrt{\dfrac{48}{\pi}}$ times the rms speed of G. If $Z_1$ represents the number of collisions made by one molecule with other molecules per unit time and the ratio of $Z_1$ for CH$_4$ and $Z_1$ for G is $x:1$, the value of `x' is

\partsec{3}{MATHEMATICS}
\sect{SECTION-I}
\qn{1} If two real and distinct roots of $(x-2019)(x-2022)+\lambda(x-2020)(x-2021)=0$ are $x_1$ and $x_2$ for some $\lambda\in\mathbb R^+$, then which of the following option(s) is/are correct -
\begin{qoptions}
\opt{A}{2020 will lie in $(x_1,x_2)$}
\opt{B}{$x_1,x_2\in(2019,2021)$}
\opt{C}{$x_1,x_2\in(2019,2022)$}
\opt{D}{$x_1,x_2\in(2019,2020)$}
\end{qoptions}

\qn{2} A straight line $y=m(x-a)$, $m\neq0$ meets a parabola $y^2=4ax$ at two points A and B. Then which of the following option(s) is/are correct
\begin{qoptions}
\opt{A}{Length AB = 4a when $m=1$}
\opt{B}{Length AB = 8a when $m=1$}
\opt{C}{Angle between normals drawn at A and B is $45^\circ$ when $m=1$}
\opt{D}{Angle between normals drawn at A and B is $90^\circ$ when $m=2$}
\end{qoptions}

\qn{3} Let $N$ be the number of ordered pairs of non-empty sets $A$ and $B$. If $A$ and $B$ satisfy

(i) $A\cup B=\{1,2,3,\ldots.12\}$

(ii) $A\cap B=\phi$

(iii) (number of elements of $A$) $\notin A$

(iv) (number of elements of $B$) $\notin B$

then choose correct options
\begin{qoptions}
\opt{A}{$N$ is a 3-digits number}
\opt{B}{Sum of the digits of $N$ is 11}
\opt{C}{When $N$ is divided by 10, remainder is 2}
\opt{D}{$N$ is an odd number}
\end{qoptions}

\qn{4} Let complex numbers $Z_1$ and $\dfrac{1}{\bar Z_1}$ lie on $|Z-Z_0|=4R$ and $|Z-Z_0|=16R$ respectively. If $Z_0$ satisfies $2|Z_0|^2=R^2+2$ and $|Z_1|^2=\dfrac pq$ where $p,q\in N$ and H.C.F of $p$ and $q$ is 1, then choose the correct option(s)
\begin{qoptions}
\opt{A}{$p>q$}
\opt{B}{$p<q$}
\opt{C}{$p$ is a prime}
\opt{D}{$q$ is a prime}
\end{qoptions}

\qn{5} If $L_1:\dfrac x1=\dfrac y2=\dfrac z3$, $P_1:x+y+z=0$, $L_2:\dfrac xa=\dfrac yb=\dfrac zc$ and line $L_2$ is projection of line $L_1$ in plane $P_1$, then which of the following option(s) is/are correct
\begin{qoptions}
\opt{A}{$a+c=0$}
\opt{B}{$a+b=0$}
\opt{C}{$b+c=0$}
\opt{D}{$a+b+c=0$}
\end{qoptions}

\qn{6} Let $P(x)$ be a monic real polynomial of degree 3. $P(x)$ has maximum at $x=0$ and $P'(x)$ has minimum at $x=1$. If $P(1)=0$, then choose correct option(s)
\begin{qoptions}
\opt{A}{$P(2)$ is negative}
\opt{B}{$P(-2)$ is negative}
\opt{C}{$P(2)$ is positive}
\opt{D}{$P(-2)$ is positive}
\end{qoptions}

\sect{SECTION-II}
\qn{1} Let $A$, $B$ and $C$ be 3 square matrices of order $3\times3$. If $A^3-6A^2+7A+B=0$ and $(\operatorname{adj}(\operatorname{adj}B))=C$, where $A=\begin{bmatrix}1&0&2\\0&2&1\\2&0&3\end{bmatrix}$, then $\left(\dfrac{\text{sum of the elements of }C}{10}\right)$ is

\qn{2} If the value of $\displaystyle\int_0^1x\left\{\dfrac{d^2}{dx^2}\left(1-x^2\right)^4\right\}dx$ is $N$, then the value of $\dfrac N2$ is

\qn{3} Let $\vec a=\hat i+2\hat j+3\hat k$, $\vec b=3\hat i+2\hat j+\hat k$ and $\vec c=\hat i+\hat j+\hat k$ be three vectors. If $\vec r\times\vec b=\vec c\times\vec b$ and $\vec r\cdot\vec a=0$, then the value of $|\vec r|^2$ is

\qn{4} Let A, B and C be the vertices of an equilateral triangle. Equation of side BC is $x+y=0$ and co-ordiantes of A are $(\alpha,\beta)$. If exterior angle bisectors of angles at B and C are $y=\left(2-\sqrt3\right)x$ and $y-1=(x+1)\left(2+\sqrt3\right)$ respectively, then the value of $(\alpha\beta)$ is

\qn{5} Two players A and B play a seires of matches. In each match it is 2 to 1 in favour of the winner of the previous match. If A wins the first match, then the probability that A will win 3 matches of the next 4 matches, is - (Assume no match ends in draw)

\qn{6} Let $T_n$ be the $n^{\text{th}}$ term of a series for $n\ge1$ and $T_1=1$, $T_{n+1}=2T_n+4^n$, $n\ge1$. If $S_n$ = Sum of first $n$ terms of the series and $S_n=\dfrac{a.4^n+b.2^n+1}{3}$ where $a,b\in I$, then the value of $\left|\dfrac ba\right|$ is

\sect{SECTION-III}
\qn{1} Let $x_0,y_0,z_0\in\mathbb R-\{0\}$. If $(x_0,y_0,z_0)$ be the solution of system of equations

$x+y-2z=0$

$2x-3y+z=0$

$x-5y+4z=k$, $k\in\mathbb R$

then the value of $\dfrac{(x_0y_0+y_0z_0+z_0x_0)}{4\left(x_0^2+y_0^2+z_0^2\right)}$ is

\qn{2} Let $f(x)$ be a polynomial function of degree $n$ such that $2+2f(x)=f(x-1)+f(x+1)$, $\forall\ x\in\mathbb R$. If $S$ = sum of all possible values of $n$, then the value of $S$ is

\qn{3} $ax^2+2hxy+by^2+2gx+2fy=0$ is a hyperbola whose center is at $(1,2)$ and the asymptotes of this hyperbola are parallel to lines $2x+3y=0$ and $x+2y=1$. If hyperbola passes through $(2,4)$, then the value of $\left[\dfrac fg\right]$ is (where [.] represents greatest integer function)

\qn{4} Let $f:[0,\pi]\to[0,\pi]$ and $g:[0,\pi]\to[0,\pi]$ be two functions defined by $f(x)=x+\sin x$ and $g$ is inverse function of $f$. If $A$ is the area of the region bounded by $y=f(x)$ and $y=g(x)$, then the value of $A$ is

\qn{5} If the sum of the series $\sec^2x+\dfrac{1}{2^2}\sec^2\left(\dfrac x2\right)+\dfrac{1}{2^4}\sec^2\left(\dfrac x{2^2}\right)+\ldots+\dfrac{1}{2^{2n-2}}\sec^2\left(\dfrac{x}{2^{n-1}}\right)$ is equal to $a^2\csc^2ax-\dfrac{4}{4^n}\csc^2\dfrac{x}{2^{n-1}}$, $a\in\mathbb N$, then the value of `a' is

\qn{6} If $f:\mathbb R-\{0,1\}\to\mathbb R$ be a function satisfying $f(x)+f\left(1-\dfrac1x\right)=\tan^{-1}x$ and $N=f\left(\dfrac12\right)$, then the value of $\left(\dfrac{8N}{\pi}\right)$ is
\end{multicols}
\end{document}
